\chapter{Engineering decision trees}
\label{app:decision-trees}

This appendix is the wall chart. It collects the universal reasoning frameworks
used throughout the book, independent of any interview question. Open it when
you need the pattern without the surrounding prose. Interview pressure maps are
in \Cref{app:interview-frameworks}; the drill philosophy is in
\Cref{app:interview-review}.

Debugging and qualification are the same philosophy at different times.
Debugging asks which margin was exhausted. Qualification asks how much margin
remains after the expected stresses. Both end on an engineering decision and a
recurrence control.

\section{Universal debugging tree}
\label{sec:tree-debugging}

\begin{dectree}
Problem
  |
Scope (unit -> lot -> vendor -> fleet)
  |
Population / time behavior
  |
Power changed?
  |-- YES --> Power ledger --> optical path
  |-- NO  --> Quality / receiver
  |
Isolation (highest-value measurement)
  |
Root cause (surviving hypothesis)
  |
Decision (contain / retune / derate / RMA / monitor)
  |
Recurrence control (ATP / SPC / telemetry / FA)
\end{dectree}

Use when BER rises, a link flaps, a lane is weak, or a corner fails. Name the
ledger before naming a component (\Cref{sec:debug-fork,sec:validation-fork}).

\section{Universal qualification tree}
\label{sec:tree-qualification}

\begin{dectree}
Requirements
  |
Bring-up
  |
Nominal characterization
  |
Margin characterization
  |
Interoperability
  |
Environmental and reliability qualification
  |
Manufacturing and ATP readiness
  |
Controlled pilot
  |
Fleet deployment and monitoring
  |
Ship / restrict / reject
\end{dectree}

Each gate removes a different uncertainty. Do not treat a bring-up pass as
margin evidence, or an HTOL pass as production readiness. Shorter ladders in
\Cref{tab:ladder,tab:ladder-gates} are grouped views of this same order
(\Cref{sec:validation-workflow,sec:interview-validation-ladder,sec:gr468}).

\section{Power-versus-quality fork}
\label{sec:tree-power-fork}

\begin{dectree}
BER or link degrade
  |
Received power changed?
  |-- YES --> Power ledger
  |           laser / coupling / connector / fiber / MUX / monitor
  |-- NO  --> Signal quality
              eye / noise / jitter / bias / EQ / RIN / Rx
  |
Highest-value measurement
  |
Decision + recurrence control
\end{dectree}

One check before retuning equalizers or bias tables
(\Cref{sec:debug-fork,sec:validation-fork}).

\section{Scope and population}
\label{sec:tree-scope-population}

\begin{dectree}
Initial symptom
  |
Scope analysis (how large?)
  |-- unit / lane / rack / site / PN / vendor / lot / fleet
  |
Technical isolation
  |
Correlation analysis (which cohort?)
  |-- lot / date code / FW / cal / platform / location
  |
Containment and corrective action
\end{dectree}

Scope sets severity and priors. Correlation after isolation unlocks contain,
pause, replace, or supplier escalate (\Cref{sec:fleet-triage,sec:yield-analysis}).

\section{Sudden versus gradual}
\label{sec:tree-time-behavior}

\begin{dectree}
Failure onset
  |
Sudden?
  |-- config / handling / firmware / mechanical / power event
  |
Gradual?
  |-- aging / drift / margin erosion / contamination / cal movement
  |
Prioritize measurements (priors, not conclusions)
\end{dectree}

Sudden and gradual are priors that reorder the bench, not root-cause claims
(\Cref{ch:failure-modes}).

\section{Transmitter, channel, or receiver}
\label{sec:tree-tx-channel-rx}

\begin{dectree}
Power or quality path chosen
  |
Golden swap / loopback
  |-- Tx --> eye / LIV / bias / wavelength / RIN
  |-- Channel --> connector / fiber / MUX / ORL
  |-- Rx --> sensitivity / TIA / EQ / host
  |
Evidence
  |
Owner + decision
\end{dectree}

Bisect domains before opening packages (\Cref{sec:debug,sec:fm-ber-increase}).

\section{Supplier qualification}
\label{sec:tree-supplier}

\begin{dectree}
Requirements (customer-visible)
  |
Characterization (distributions, not samples)
  |
Margin under stress
  |
Qualification (named mechanisms)
  |
Production readiness (FAIR, ATP, SPC)
  |
Fleet monitoring (RMA, telemetry)
\end{dectree}

Customer view measures external behavior. Vendor view owns internals
(\Cref{sec:interview-customer-vendor,sec:supplier-exec}).

\section{Supplier escape and containment}
\label{sec:tree-escape}

\begin{dectree}
Escape detected
  |
Contain scoped population
  |
Scope (unit / lot / vendor / site)
  |
Evidence (plane, corner, telemetry)
  |
Root cause class
  |
Supplier / FA path
  |
ATP or screen update
  |
SPC / reaction plan
  |
Fleet monitoring
\end{dectree}

Contain first when the population can grow. The system owner keeps
responsibility for evidence quality and verifying corrective action
(\Cref{sec:escaped-defects,sec:fleet-triage}).

\section{Margin-budget flow}
\label{sec:tree-margin-budget}

\begin{dectree}
Nominal system margin
  |
Temperature debit
  |
Voltage / power-quality debit
  |
Channel / connector debit
  |
Manufacturing variation
  |
Aging / wear
  |
Interoperability variation
  |
Remaining margin
  |
Above deployment requirement?
  |-- YES --> proceed
  |-- NO  --> redesign / restrict / recalibrate / reject
\end{dectree}

Design allocates budgets. Validation often measures the net externally visible
result. Do not double-count internal penalties the test cannot separate
(\Cref{sec:laser-margin-erosion,sec:link-budget}).

\section{Black-box versus engineering access}
\label{sec:tree-blackbox-access}

\begin{dectree}
Bookended product
  |
End-to-end qualification
  |-- BER / FEC / telemetry / sensitivity
  |-- environment / interoperability
  |
Enough confidence to decide?
  |-- YES --> deployment decision
  |-- NO  --> request engineering access
              Tx-only / Rx-only / breakout / diagnostics
              |
              Isolate margin
\end{dectree}

An optical eye is measured externally with suitable access. Do not assume the
module reports a conventional eye unless that capability exists
(\Cref{sec:interview-customer-vendor}).

\section{Recurrence-control loop}
\label{sec:tree-recurrence}

\begin{dectree}
Production escape
  |
Contain risk now
  |
Investigate mechanism
  |
Correct process
  |
Update screening or ATP
  |
Verify next lot
  |
Monitor fleet
\end{dectree}

Containment, root cause, and prevention are three different actions
(\Cref{sec:escaped-defects}).

\section{Production feedback loop}
\label{sec:tree-production-loop}

\begin{dectree}
Design requirements
  |
Qualification
  |
ATP
  |
Production data
  |
Fleet data
  |
Failure analysis
  |
Updated limits or screens
  |
Next production cycle
\end{dectree}

Production validation is replayable and decision-oriented
(\Cref{sec:supplier-exec,sec:hvm-test}).

\section{Measurement-selection loop}
\label{sec:tree-measurement-loop}

\begin{dectree}
Current evidence
  |
Rank hypotheses
  |
Choose measurement that removes
the largest number of hypotheses
  |
Update beliefs
  |
Enough evidence to decide?
  |-- NO  --> next measurement
  |-- YES --> take action
\end{dectree}

\keyidea{Great engineers perform the minimum measurement that eliminates the
largest number of hypotheses.}

\section{Unknown failure}
\label{sec:tree-unknown}

\begin{dectree}
Observation (facts, not hopes)
  |
Hypotheses (short ranked list)
  |
Highest-value measurement
  |
Evidence (what died)
  |
Decision with today's confidence
  |
Continue measuring?
  |-- YES --> next measurement
  |-- NO  --> control + owner + residual risk
\end{dectree}

Unknown mechanism is not a freeze. Decide with weighted evidence
(\Cref{sec:interview-under-uncertainty,sec:tree-measurement-loop}).

\section{Before you start: three checklists}
\label{sec:tree-checklists}

\engcheck{Before debugging}{%
Scope $\cdot$
unit/rack/lot/vendor/fleet $\cdot$
sudden or gradual $\cdot$
power changed? $\cdot$
signal-quality uncertainty $\cdot$
highest-value measurement $\cdot$
decision unlocked $\cdot$
containment now $\cdot$
recurrence control.}

\engcheck{Before qualification}{%
Nominal function $\cdot$
operating margin $\cdot$
environmental stresses $\cdot$
interoperability $\cdot$
aging/reliability $\cdot$
manufacturing variation $\cdot$
ATP escape detection $\cdot$
controlled pilot $\cdot$
fleet telemetry.}

\engcheck{Before production}{%
Requirements met $\cdot$
margin demonstrated $\cdot$
environment covered $\cdot$
interop verified $\cdot$
reliability risk addressed $\cdot$
lot variation characterized $\cdot$
ATP defined $\cdot$
supplier controls reviewed $\cdot$
pilot exit criteria $\cdot$
fleet monitoring ready.}

\keyidea{These trees are the book's reusable method. Chapters supply the
physics and the numbers. This appendix supplies the order of thought.}
